\chapter*{引言}
\addcontentsline{toc}{chapter}{引言}
\markboth{引言}{引言}

\section{平坦连接、D-模与晶体}

\ob{0.1.1}{n:0.1.1}设 $M$ 是光滑流形，$V$ 是其上的向量丛。回忆：$V$ 上的平坦连接是满足 Leibniz 法则、且曲率 $[\nabla,\nabla]=0$ 的映射
\[
\nabla:V\to V\otimes\Omega^1_M .
\]
把连接映射对偶化，得到映射
\[
T_M\otimes V\to V .
\]
连接的平坦性保证 $V$ 由此成为向量场 Lie 代数上的模。等价地，$V$ 成为 $M$ 上微分算子代数 $\Diff_M$ 上的模。

这一概念立刻推广到特征零的光滑代数簇。在这样的簇上，D-模定义为微分算子层上的模，且作为 $\mathcal{O}$-模是拟凝聚的。于是代数簇上 D-模的概念推广了带平坦连接的向量丛，并编码了系数为多项式的线性微分方程组的资料。光滑代数簇上 D-模的研究是一个内容极其丰富的理论，在表示论等诸多领域有应用；带平坦连接的向量丛的微分几何中的许多想法，都可以搬到这一框架里。

然而，上述 D-模的进路有几处困难。例如，除光滑概形外，还需要考虑奇异概形上带平坦连接的层。微分算子代数在奇异簇上仍有定义，但它的模范畴并不是我们想要的那个范畴（比如，这个代数一般不是 Noether 的）。另一方面，即使对光滑代数簇，如何对非线性对象——例如层的范畴——定义连接，也并不清楚。

\ob{0.1.2}{n:0.1.2}\textbf{平行移动。}更好的定义来自微分几何中平坦连接的另一种解读，即平行移动。

设 $M$ 是光滑流形，其上带平坦连接的向量丛 $V$，并取道路 $\gamma:[0,1]\to M$。于是得到纤维之间的同构
\[
\Pi_\gamma:V_{\gamma(0)}\simeq V_{\gamma(1)},
\]
它只依赖道路的同伦类。这一构造可以换个说法。取 $M$ 内的小球 $B$。平行移动同构只依赖道路的同伦类，而 $B$ 可缩，所以沿 $B$ 内的道路，得到纤维的一个相容等同
\[
V_x\simeq V_y\qquad (x,y\in B).
\]
粗略地说，连接的数据给出了流形上「邻近」点处纤维的等同。

沿着这一想法，Grothendieck [Gr] 给出了平行移动这一概念的纯代数类比，并且本质地用到概形理论（而不只是簇）：他引入了概形 $X$ 的 $R$-点之间的无穷小接近关系。具体地说，两个 $R$-点 $x,y:\Spec(R)\to X$ 称为无穷小接近的，如果它们限制到 $\Spec(\mathrm{red}R)$ 上相同，其中 $\mathrm{red}R$ 是 $R$ 模去幂零根所得的商环。

概形 $X$ 上的晶体，按定义就是 $X$ 上关于无穷小接近关系等变的拟凝聚层。更确切地说，$X$ 上的晶体是一个拟凝聚层 $\mathcal{F}$，附带同构
\[
x^*(F)\simeq y^*(F)
\]
的数据，对任意两个无穷小接近的点 $x,y:\Spec(R)\to X$ 给出，并要求满足余循环条件。

Grothendieck 证明：在光滑代数簇上，晶体的 Abel 范畴等价于微分算子环上左模范畴。就此而言，晶体给出了带平坦连接的层的一个更基本的定义。

晶体范畴的一个突出特点是，Kashiwara 引理已经内置于它的定义之中：对概形的闭浸入 $i:Z\to X$，$Z$ 上晶体的范畴等价于 $X$ 上作为集合支撑在 $Z$ 上的晶体的范畴。这一观察使我们能把概形上晶体的研究化归到光滑概形的情形：把给定的概形（局部地）浸入一个光滑概形即可。

\ob{0.1.3}{n:0.1.3}本文在导出代数几何的框架里发展晶体的理论。那里考虑的不是通常的环，而是导出环，即 $E_\infty$ 环谱。由于在特征零的域 $k$ 上工作，我们取连接的交换 DG $k$-代数作为导出环的模型（相应地，使用术语「DG 概形」而不是「导出概形」）。关键的想法是：应当把导出环的高阶同伦群看作幂零元概念的推广。

于是，仿照 Simpson [Si]，对 DG 概形 $X$，定义它的 de Rham 预叠 $X_{\dR}$ 为导出环范畴上的函子
\[
X_{\dR}:R\mapsto X(\clred R),
\]
其中
\[
\clred R:=\mathrm{red}\big(\pi_0(R)\big)
\]
是与 $R$ 的底层经典环相对应的既约环。换言之，按 [GL:Stacks] 的术语，$X_{\dR}$ 是一个预叠。

我们定义 $X$ 上的晶体为预叠 $X_{\dR}$ 上的拟凝聚层。预叠上拟凝聚层的理论见 [Lu1, Sect.~2]，简要的回顾见 [GL:QCoh, Sect.~1.1]。

上述定义与前面提到的 Grothendieck 的定义并不一致：后者把映射 $\Spec(R)\to X$ 指定到等价类为止，前者则只指定映射 $\Spec(\clred R)\to X$。不过我们将证明：当 $X$ 最终余连接，即其结构环只有有限多个非零同伦群时，晶体的这两种定义等价。\footnote{当 $X$ 不是最终余连接时，这两种概念不同，正确的一种是经由 $X_{\dR}$ 定义的那一种。}

\ob{0.1.4}{n:0.1.4}晶体范畴虽然与 D-模范畴等价，却提供了一个更灵活的发展理论的框架。定义立刻推广到非光滑概形，相应的范畴性质良好（例如，任意概形上晶体的范畴是局部 Noether 的）。

\begin{annote}
为什么需要导出环？经典概形的 de Rham 预叠只看见既约部分，单看点集并不需要高阶同伦群。导出环的用处在于：沿非平坦映射取纤维积、取形式完备化、把晶体沿任意态射拉回时，结构环会出现幂零元；在导出代数几何中，这些幂零元自然由高阶同伦群承载。于是 connective DG 代数既保留了经典情形，又在取极限时封闭。\hfill【来源：原书 §0.1.3 及脚注 1；[GR1, Sect.~6.1]（A）】
\end{annote}


设 $f:X\to Y$ 是 DG 概形的映射。我们将构造自然的拉回函子
\[
f^\dagger:\Crys(Y)\to\Crys(X).
\]
事实上，我们将把 $X\mapsto\Crys(X)$ 延拓为从 $\DGSch^{\mathrm{op}}$ 到稳定 $\infty$-范畴的函子。后者使我们不只在 DG 概形上，而且在任意预叠上定义晶体。

此外，晶体的概念立刻推广到非线性与范畴化的情形：我们同样可以定义 $X$ 上的概形晶体或范畴晶体。

\section{左晶体与右晶体}

\ob{0.2.1}{n:0.2.1}回忆：在光滑代数簇 $X$ 上，除通常的（即左）D-模外，还可以考虑右 D-模的范畴。这两个范畴等价，相应的函子是沿函数环与对偶化线丛 $\omega_X$ 作张量积。然而这个等价不保持到拟凝聚层的遗忘函子。因此，我们可以考虑一个抽象的 D-模范畴，它带有两个不同的实现函子，分别到拟凝聚层。在左实现中，D-模的拉回函子成为拟凝聚层上的 $*$-拉回函子；在右实现中，它成为 $!$-拉回函子。

事实证明，「右」实现比「左」实现有几个优点。主要的优点也许是：「右」实现给 D-模范畴赋予了一个函子性非常好的 $t$-结构。特别地，在 Riemann--Hilbert 对应下，这个 $t$-结构成为反常 $t$-结构。

\ob{0.2.2}{n:0.2.2}于是可以问：在任意 DG 概形上，是否同样有「左」「右」晶体？答案是：两个范畴确实都可以非常一般地定义。

左晶体就是 §0.1.3 中定义的对象。但为了定义右晶体，需要把通常的拟凝聚层范畴换成它的重整化版本，即 [IndCoh] 中引入的归纳凝聚层范畴。

归纳凝聚层范畴 $\IndCoh(X)$ 对（几乎）局部有限型的（导出）概形性质良好，所以右晶体只在这样的 DG 概形上定义，进而只在此性质的预叠上定义。

回忆 [IndCoh, Sect.~5]：对 DG 概形之间的映射 $f:X\to Y$，有 $!$-拉回函子
\[
f^!:\IndCoh(Y)\to\IndCoh(X).
\]
附着 $X\mapsto\IndCoh(X)$ 是从 $\DGSch^{\mathrm{op}}$ 到稳定 $\infty$-范畴的函子，因而可以延拓为从预叠范畴出发的函子。

对 DG 概形 $X$，定义右晶体的范畴 $\Crys^r(X)$ 为 $\IndCoh(X_{\dR})$。也可以按 Grothendieck 的方式重新表述这个定义：$X$ 上的右晶体是对象 $F\in\IndCoh(X)$，附带对每对无穷小接近的点 $x,y:\Spec(R)\to X$ 的同构
\begin{equation}
x^!(F)\simeq y^!(F),\tag{0.1}
\end{equation}
并要求满足（$\infty$-范畴版本的）余循环条件。可以证明，与左晶体的情形不同，这确实给出了右晶体的等价定义，而不需要任何余连接假设。

\begin{annote}
为什么右晶体必须用归纳凝聚层？因为条件 (0.1) 要求对任意 $x,y$ 比较 $x^!(F)$ 与 $y^!(F)$，而 $!$-拉回一般只在归纳凝聚层上定义；$\QCoh$ 上的 $*$-拉回无法表达「支撑在无穷小邻域上」的信息。另一方面，$\Upsilon_X:\QCoh(X)\to\IndCoh(X)$ 在 $X$ 局部几乎有限型时把两种实现联系起来，故左右晶体最终仍然等价（§2.4）。\hfill【来源：原书 §0.2.2--0.2.3；[IndCoh, Sect.~4--5]（A）】
\end{annote}


\ob{0.2.3}{n:0.2.3}右晶体的范畴既已定义，就可以问它是否与左晶体的范畴等价。答案同样是肯定的。具体地说，对几乎有限型的任意 DG 概形 $X$，与对偶化复形 $\omega_X$ 作张量积给出函子
\[
\Upsilon_X:\QCoh(X)\to\IndCoh(X),
\]
它把拟凝聚层上的 $*$-拉回与归纳凝聚层上的 $!$-拉回联系起来。

虽然对单个 $S$，函子 $\Upsilon_S$ 一般不是等价（除非 $S$ 光滑），但对所有映入 $X$ 的 de Rham 预叠的 DG 概形取这些映射的整体，得到的却是左晶体与右晶体之间的等价。

于是，正如光滑簇的情形，可以认为对每个 DG 概形 $X$ 都附着一个范畴 $\Crys(X)$，它带两个「实现」函子
\[
\begin{tikzcd}
 & \Crys(X) \arrow[dl, "\oblv^l_X"'] \arrow[dr, "\oblv^r_X"] & \\
 \QCoh(X) \arrow[rr, "\Upsilon_X"'] & & \IndCoh(X)
\end{tikzcd}
\]
然而在非光滑概形的情形，「右」实现所对应的 $\Crys(X)$ 上 $t$-结构的优点更加明显。

\ob{0.2.4}{n:0.2.4}\textbf{历史注。}据我们所知，经由右晶体研究 D-模的进路最早由 A. Beilinson 在 90 年代初提出，当时是在 Abel 范畴的层面。

此后一段时间里，人们误以为不能用左晶体来定义 D-模，理由是 $t$-结构不相容。然而 J. Lurie 说明：只要放弃 $t$-结构、直接定义相应的稳定 $\infty$-范畴，左晶体同样适用。

\section{晶体/D-模的理论}

让我们说明本文及其续篇 [GR2] 所发展的理论的整体结构。

\ob{0.3.1}{n:0.3.1}对每个（局部几乎有限型的）预叠 $Y$，我们赋予一个稳定 $\infty$-范畴
\[
Y\rightsquigarrow\Crys(Y).
\]
这个范畴有两个实现函子：到 $\QCoh(Y)$ 的左实现函子，与到 $\IndCoh(Y)$ 的右实现函子。二者由如下交换图联系：
\[
\begin{tikzcd}
 & \Crys(Y) \arrow[dl, "\oblv^l_Y"'] \arrow[dr, "\oblv^r_Y"] & \\
 \QCoh(Y) \arrow[rr, "\Upsilon_Y"'] & & \IndCoh(Y)
\end{tikzcd}
\]
其中 $\Upsilon_Y$ 是与对偶化复形 $\omega_Y$ 作张量积给出的函子 $\QCoh(Y)\to\IndCoh(Y)$。

\ob{0.3.2}{n:0.3.2}把 $\Crys(Y)$ 附着到 $Y$ 上的做法在多种意义下具有函子性。对映射 $f:Y_1\to Y_2$，有拉回函子
\[
f^\dagger:\Crys(Y_2)\to\Crys(Y_1),
\]
它对 $f$ 具有函子性；也就是说，这一附着给出函子
\[
\Crys^\dagger_{\PreStk}:(\PreStk)^{\mathrm{op}}\to\DGCat_{\mathrm{cont}}.
\]
D-模上的拉回函子与实现函子相容。具体地说，有交换图
\[
\begin{tikzcd}
 \Crys(Y_1) & \Crys(Y_2) \arrow[l, "f^\dagger"'] \\
 \QCoh(Y_1) \arrow[u, "\oblv^l_{Y_1}"] & \QCoh(Y_2) \arrow[l, "f^*"] \arrow[u, "\oblv^l_{Y_2}"']
\end{tikzcd}
\]
以及
\[
\begin{tikzcd}
 \Crys(Y_1) & \Crys(Y_2) \arrow[l, "f^\dagger"'] \\
 \IndCoh(Y_1) \arrow[u, "\oblv^r_{Y_1}"] & \IndCoh(Y_2) \arrow[l, "f^!"] \arrow[u, "\oblv^r_{Y_2}"']
\end{tikzcd}
\]
此外，这种相容性本身对 $f$ 具有函子性；也就是说，有一个自然交换的函子图
\[
\begin{tikzcd}
 & \Crys^\dagger_{\PreStk} \arrow[dl, "\oblv^l"'] \arrow[dr, "\oblv^r"] & \\
 \QCoh_{\PreStk} \arrow[rr, "\Upsilon"'] & & \IndCoh^!_{\PreStk}
\end{tikzcd}
\]

\ob{0.3.3}{n:0.3.3}上述这部分理论在本文中构造。也就是说，本文关心的是附着
\[
Y\rightsquigarrow\Crys(Y)
\]
以及拉回运算。因此，本文发展的是晶体/D-模的局部理论。

不过，除函子 $f^\dagger$ 之外，我们还期待有推前函子 $f_{\dR,*}$，并且两者须满足各种相容关系。后者将在 [GR2] 中完成。这里先指出这套合并理论的主要构件。

\ob{0.3.4}{n:0.3.4}对预叠之间概形性的拟紧映射 $f:Y_1\to Y_2$，有 de Rham 推前函子
\[
f_{\dR,*}:\Crys(Y_1)\to\Crys(Y_2),
\]
它对 $f$ 具有函子性。这一附着给出另一个函子
\[
\big(\Crys_{\dR,*}\big)_{\PreStk^{\mathrm{sch\text{-}qc}}}:\PreStk^{\mathrm{sch\text{-}qc}}\to\DGCat_{\mathrm{cont}},
\]
其中 $\PreStk^{\mathrm{sch\text{-}qc}}$ 是 $\PreStk$ 的非全子范畴，其 1-态射限制为概形性的拟紧映射。

设 $Y=X$ 是 DG 概形。\footnote{更一般地，可以取 $Y$ 为具有形变理论的预栈。}此时遗忘函子
\[
\oblv^r_Y:\Crys(Y)\to\IndCoh(Y)
\]
有左伴随，记作
\[
\indf^r_Y:\IndCoh(Y)\to\Crys(Y),
\]
称为归纳函子。

归纳函子与 de Rham 推前相容。具体地说，有交换图
\[
\begin{tikzcd}
 \IndCoh(Y_1) \arrow[r, "f^{\IndCoh}_*"] \arrow[d, "\indf^r_{Y_1}"'] & \IndCoh(Y_2) \arrow[d, "\indf^r_{Y_2}"] \\
 \Crys(Y_1) \arrow[r, "f_{\dR,*}"'] & \Crys(Y_2)
\end{tikzcd}
\]
这种相容性本身也具有函子性；也就是说，有函子的自然变换
\[
\indf^r:(\IndCoh_*)_{\PreStk^{\mathrm{sch\text{-}qc}}}\to\big(\Crys_{\dR,*}\big)_{\PreStk^{\mathrm{sch\text{-}qc}}}.
\]

\ob{0.3.5}{n:0.3.5}当 $f$ 固有时，函子 $(f_{\dR,*},f^\dagger)$ 构成伴随对；当 $f$ 光滑时，若 $n$ 是 $f$ 的相对维数，则 $(f^\dagger[-2n],f_{\dR,*})$ 构成伴随对。

一般地，这两个函子并不伴随，但满足基变换公式。正如 [IndCoh, Sect.~5.1] 所解释的，编码基变换公式的函子性的一种方式是考虑对应范畴。具体地说，令 $(\PreStk)^{\mathrm{corr:all,sch\text{-}qc}}$ 是这样的 $\infty$-范畴：其对象是局部有限型的预叠，从 $Y_1$ 到 $Y_2$ 的态射由对应
\[
\begin{tikzcd}
 & Z \arrow[dl, "f"'] \arrow[dr, "g"] & \\
 Y_2 & & Y_1
\end{tikzcd}
\]
给出，其中 $f$ 概形性且拟紧，$g$ 任意。这个范畴中的复合由对应的纤维积给出。于是，函子 $\Crys^\dagger$ 与 $\Crys_{\dR,*}$ 的相容基变换公式是一个函子
\[
\Crys_{(\PreStk)^{\mathrm{corr:all,sch\text{-}qc}}}:(\PreStk)^{\mathrm{corr:all,sch\text{-}qc}}\to\DGCat_{\mathrm{cont}},
\]
连同它在 $(\PreStk)^{\mathrm{op}}$ 上的限制与 $\Crys^\dagger_{\PreStk}$ 的等同，以及在 $\PreStk^{\mathrm{sch\text{-}qc}}$ 上的限制与 $(\Crys_{\dR,*})_{\PreStk^{\mathrm{sch\text{-}qc}}}$ 的等同。

\section{扭化}

\ob{0.4.1}{n:0.4.1}除 D-模外，考虑扭 D-模往往也很重要。例如在表示论中，Beilinson 与 Bernstein 的局部化定理把具有固定中心特征 $\chi$ 的约化 Lie 代数 $\mathfrak{g}$ 的表示范畴，等同于旗簇 $G/B$ 上扭 D-模的范畴，其扭化由 $\chi$ 决定。

在光滑簇的情形，扭化与扭 D-模的理论由 Beilinson 与 Bernstein [BB] 引入。扭化的重要例子来自线丛的复张量幂。对光滑簇 $X$，扭化构成一个 Picard 群胚，它可以这样描述。令 $T$ 为次数 1 与 2 上的层复形
\[
T:=\Omega^1\to\Omega^{2,\mathrm{cl}},
\]
其中 $\Omega^1$ 是 $X$ 上 1-形式层，$\Omega^{2,\mathrm{cl}}$ 是闭 2-形式层，映射是 de Rham 微分。于是扭化的 Picard 群胚的对象空间由 $H^2(X,T)$ 给出；对一个给定对象，同构空间是 $H^1(X,T)$。

\ob{0.4.2}{n:0.4.2}本文最后两节在导出（因而特别是非光滑）的框架里发展扭化与扭晶体的理论。我们给出扭化概念的几种等价表述，并证明它们在光滑簇的情形与 [BB] 中的定义一致。

对（几乎局部有限型的）预叠 $Y$，定义扭化为 de Rham 预叠 $Y_{\dR}$ 上的 $\mathbb{G}_m$-gerbe，并附带其到 $Y$ 的拉回的平凡化。$Y$ 上的线丛 $L$ 给出一个扭化：$Y_{\dR}$ 上的 gerbe 平凡，而 $Y$ 上的平凡化由 $L$ 给出。

给定扭化 $T$，$Y$ 上 $T$-扭晶体的范畴定义为 $Y_{\dR}$ 上被 $T$ 所给的 $\mathbb{G}_m$-gerbe 扭化的（归纳凝聚或拟凝聚）层范畴。

\section{内容概览}

下面逐节说明本文的内容。

\ob{0.5.1}{n:0.5.1}在 §1 中，对预叠 $Y$ 定义 de Rham 预叠 $Y_{\dR}$，并建立它的一些基本性质。最重要的是证明：若 $Y$ 局部几乎有限型，则 $Y_{\dR}$ 也局部几乎有限型。

\ob{0.5.2}{n:0.5.2}在 §2 中，把左晶体定义为 de Rham 预叠上的拟凝聚层；在局部几乎有限型的情形，把右晶体定义为 de Rham 预叠上的归纳凝聚层。后者之所以有意义，是因为 §1 证明了：局部几乎有限型的预叠，其 de Rham 预叠也局部几乎有限型。在这一情形，我们证明左晶体与右晶体的范畴等价。此外还证明这一框架下的 Kashiwara 引理。

\ob{0.5.3}{n:0.5.3}在 §3 中，证明晶体范畴满足 h-下降（特别地，fppf 下降）。我们还引入预叠 $Y$ 的无穷小群胚，即自然映射 $Y\to Y_{\dR}$ 的 Čech 神经。具体地说，$Y$ 的无穷小群胚由
\[
(Y\times Y)^\wedge_Y\rightrightarrows Y
\]
给出，其中 $(Y\times Y)^\wedge_Y$ 是 $Y\times Y$ 沿对角线的形式完备化。

§3 的大部分内容专门讨论 $Y$ 是 ind-概形的情形。无穷小群胚上的层，就是 $Y$ 上关于无穷小接近等价关系等变的层。在归纳凝聚层的情形，这个范畴等价于右晶体。然而一般说来，无穷小群胚上的拟凝聚层并不等价于左晶体。我们证明：当 $Y$ 是最终余连接的 DG 概形或经典形式光滑的预叠时，$Y$ 的无穷小群胚上的拟凝聚层等价于左晶体。特别地，对经典概形，这个等价成立。

我们还定义从 $\QCoh(Y)$ 与 $\IndCoh(Y)$ 到 $Y$ 上晶体的归纳函子。在归纳凝聚层的情形，归纳函子是遗忘函子的左伴随，且右晶体的范畴等价于相应单子上的模范畴。对 $\QCoh$ 与左晶体，当 $Y$ 是最终余连接的 DG 概形时，类似结果成立。

\ob{0.5.4}{n:0.5.4}在 §4 中，证明晶体范畴有两个自然的 $t$-结构：一个与到 $\QCoh$ 的左实现相容，另一个与到 $\IndCoh$ 的右实现相容。对拟紧 DG 概形，这两个 $t$-结构相差一个有界幅度。

我们还证明：对仿射 DG 概形，晶体范畴等价于其关于右 $t$-结构的心上的导出范畴。

\ob{0.5.5}{n:0.5.5}在 §5 中，把 DG 概形上作用在 $\IndCoh$（相应地，$\QCoh$）上、负责右（相应地，左）晶体范畴的单子，与微分算子层联系起来。

作为推论，我们把晶体范畴与 D-模的导出范畴联系起来。

\ob{0.5.6}{n:0.5.6}在 §6 中，定义预叠 $Y$ 上扭化的 Picard 群胚为 de Rham 预叠 $Y_{\dR}$ 上、在 $Y$ 上平凡化的 $\mathbb{G}_m$-gerbe 的 Picard 群胚。然后给出这个定义的几种等价表述。例如，利用某一版本的指数映射，证明扭化的 Picard 群胚等价于 de Rham 预叠 $Y_{\dR}$ 上、在 $Y$ 上平凡化的 $\mathbb{G}_a$-gerbe 的 Picard 群胚。特别地，这使扭化自然成为一个 $k$-线性 Picard 群胚。

此外，利用扭化的 $\mathbb{G}_a$-gerbe 描述，我们把扭化的 $\infty$-群胚等同于
\[
\tau^{\le 2}\Big(H_{\dR}(Y)\times_{H(Y)}\{*\}\Big)[2],
\]
其中 $H_{\dR}(Y)$ 是 $Y$ 的 de Rham 上同调，$H(Y)$ 是 $Y$ 的凝聚上同调。特别地，对光滑经典概形，这表明这里定义的扭化概念与 [BB] 中的定义一致。

最后，我们证明：局部有限型的 DG（ind-）概形 $X$ 上扭化的范畴，等同于其无穷小群胚的中心扩张的范畴。

\ob{0.5.7}{n:0.5.7}在 §7 中，定义扭晶体的范畴并建立其基本性质。特别地，我们证明关于晶体的大部分结果都可以搬到扭化的情形。

\section{约定与记号}

我们的约定紧跟 [GR1]。这里回顾其中最必要的部分。

\ob{0.6.1}{n:0.6.1}\textbf{基域。}全文在固定的特征零基域 $k$ 上工作。

\ob{0.6.2}{n:0.6.2}\textbf{$\infty$-范畴。}所谓 $\infty$-范畴，总是指 $(\infty,1)$-范畴。我们有时把 $\infty$-范畴简称为「范畴」，这是轻微的语言滥用。我们对 $\infty$-范畴的使用不依赖模型，但心目中取的是它们作为准范畴的实现。准范畴意义下 $\infty$-范畴的基本参考书是 [Lu0]。

记 $\infty\text{-}\Grpd$ 为 $\infty$-群胚的 $\infty$-范畴，它等同于 [Lu0] 记号下的空间范畴 $\mathcal{S}$。

对 $\infty$-范畴 $\mathcal{C}$ 与 $x,y\in\mathcal{C}$，记 $\Maps_{\mathcal{C}}(x,y)\in\infty\text{-}\Grpd$ 为相应的映射空间。记 $\Hom_{\mathcal{C}}(x,y)$ 为集合 $\pi_0(\Maps_{\mathcal{C}}(x,y))$，即 [Lu0] 中记作 $\Hom_{h\mathcal{C}}(x,y)$ 的东西。

稳定 $\infty$-范畴 $\mathcal{C}$ 自然地充实于谱。此时对 $x,y\in\mathcal{C}$，记 $\Maps_{\mathcal{C}}(x,y)$ 为 $x$ 到 $y$ 的映射谱。特别地，映射空间等于该映射谱的 $\Omega^\infty$。

在固定的 $\infty$-范畴 $\mathcal{C}$ 中工作时，对两个对象 $x,y\in\mathcal{C}$，我们把 $\Maps_{\mathcal{C}}(x,y)$ 的点称为同构，而 [Lu0] 中称之为等价。也就是说，同构是有同伦逆的映射。我们保留「等价」一词，专指 $\infty$-范畴之间的（同伦）等价。

\ob{0.6.3}{n:0.6.3}\textbf{DG 范畴。}我们关于 DG 范畴的约定遵循 [IndCoh, Sect.~0.6.4]。所谓 DG 范畴，指的是 $k$ 上可表现的 DG 范畴；特别地，我们的 DG 范畴都假设余完备。除另有说明，我们只考虑 DG 范畴之间的连续函子（即与直和交换、等价地与所有余极限交换的精确函子）。换言之，我们在 [GL:DG] 记号下的范畴 $\DGCat_{\mathrm{cont}}$ 中工作。\footnote{可以把 $\DGCat_{\mathrm{cont}}$ 换成（等价的）稳定可表现 $\infty$-范畴的 $(\infty,1)$-范畴，这些范畴张量于 $\Vect$ 上，态射为保余极限的函子。}

记 $\Vect$ 为 $k$-向量空间复形的 DG 范畴。范畴 $\DGCat_{\mathrm{cont}}$ 有自然的对称单幺结构，其单位是 $\Vect$。

对带 $t$-结构的 DG 范畴 $\mathcal{C}$，记 $\mathcal{C}^{\le n}$（相应地，$\mathcal{C}^{\ge m}$、$\mathcal{C}^{\le n,\ge m}$）为 $\mathcal{C}$ 中由满足 $H^i(x)=0$（$i>n$；相应地，$i<m$；$(i>n)\wedge(i<m)$）的对象 $x$ 张成的全子范畴。包含 $\mathcal{C}^{\le n}\hookrightarrow\mathcal{C}$ 有右伴随，记作 $\tau^{\le n}$；其余类似。

有一个从 $\DGCat_{\mathrm{cont}}$ 到稳定 $\infty$-范畴与连续精确函子范畴的全忠实函子。如此得到的稳定 $\infty$-范畴充实于 $\Vect$。因此我们常把映射谱 $\Maps_{\mathcal{C}}(x,y)$ 看作 $\Vect$ 的对象；前者由后者经 Dold--Kan 对应得到。

\ob{0.6.4}{n:0.6.4}\textbf{（预）叠与 DG 概形。}我们关于（预）叠与 DG 概形的约定遵循 [GL:Stacks]：

记 $\DGSch_{\mathrm{aff}}$ 为 $k$ 上连接交换 DG 代数范畴的反范畴。

预叠的范畴 $\PreStk$ 按定义是全体函子
\[
(\DGSch_{\mathrm{aff}})^{\mathrm{op}}\to\infty\text{-}\Grpd
\]
构成的范畴。

记 $^{<\infty}\DGSch_{\mathrm{aff}}$ 为 $\DGSch_{\mathrm{aff}}$ 中由最终余连接对象张成的全子范畴。

回忆：最终余连接的仿射 DG 概形 $S=\Spec(A)$ 称为几乎有限型的，如果
\begin{itemize}[nosep]
\item $H^0(A)$ 在 $k$ 上有限型；
\item 每个 $H^i(A)$ 作为 $H^0(A)$-模有限生成。
\end{itemize}
记 $^{<\infty}\DGSch_{\mathrm{aff}}^{\mathrm{aft}}$ 为 $^{<\infty}\DGSch_{\mathrm{aff}}$ 中由几乎有限型的概形张成的全子范畴，并记 $\PreStk_{\mathrm{laft}}$ 为全体函子
\[
\big({}^{<\infty}\DGSch_{\mathrm{aff}}^{\mathrm{aft}}\big)^{\mathrm{op}}\to\infty\text{-}\Grpd
\]
构成的范畴。正如 [GL:Stacks, Sect.~1.3.11] 所解释的，$\PreStk_{\mathrm{laft}}$ 经由适当的 Kan 延拓自然地成为 $\PreStk$ 的子范畴。

为了应用 [IndCoh] 中发展的归纳凝聚层形式体系，本文大部分内容假设所考虑的预叠局部几乎有限型。遇到不如此假设的地方，我们会明确指出。

\ob{0.6.5}{n:0.6.5}\textbf{既约环。}记 $({}^{\mathrm{red}}\Sch_{\mathrm{aff}})^{\mathrm{op}}\subset(\DGSch_{\mathrm{aff}})^{\mathrm{op}}$ 为既约离散环的范畴。包含函子有自然的左伴随
\[
(-)^{\mathrm{cl,red}}:(\DGSch_{\mathrm{aff}})^{\mathrm{op}}\to({}^{\mathrm{red}}\Sch_{\mathrm{aff}})^{\mathrm{op}},
\]
由
\[
S\mapsto H^0(S)/\mathrm{nilp}(H^0(S))
\]
给出，其中 $\mathrm{nilp}(H^0(S))$ 是 $H^0(S)$ 中幂零元构成的理想。

\section{致谢}

我们感谢 Jacob Lurie 的多次有益讨论。他的想法对本论文影响如此之深，以至于很难一一指出哪些具体陈述是直接向他学到的。

D.G. 的研究由 NSF 基金 DMS-1063470 资助。
